考虑理想无像差、无损且角孔径 \(0<\alpha<\pi/2\) 的矢量光学聚焦系统。焦区Debye角谱定义为
\[
\boldsymbol E(\boldsymbol r)=C\int_0^\alpha\!d\theta\,\sin\theta
\int_0^{2\pi}\!d\varphi\ \boldsymbol e_\theta(\theta,\varphi)
e^{ik\boldsymbol r\cdot\boldsymbol s},
\]
其中 \(\boldsymbol s=(\sin\theta\cos\varphi,\sin\theta\sin\varphi,\cos\theta)\)，
\(\boldsymbol e_\theta=(\cos\theta\cos\varphi,\cos\theta\sin\varphi,-\sin\theta)\)，\(C\) 可为复常数。该幅度约定已经吸收物镜的apodization。
(1) 证明各角谱分量横向于传播方向，给出整束场满足的无源散度及Helmholtz方程。
(2) 在轴上 \(\rho=0\) 计算电场的横向分量和纵向分量，特别给出焦点 \(z=0\) 的完整解析公式。
(3) 在焦平面 \(z=0\) 用Bessel积分化为径向 \(E_r\) 和纵向 \(E_z\) 两个一重积分，说明轴向径向场严格消失但轴向纵场非零。
(4) 取 \(\alpha\ll1\)、\(k\rho\ll1\) 求近轴纵场与径向分量的主导渐近，比较径向与方位偏振入射的纵场差别，并指出本题角谱模型的偏振与几何光学适用性。
